
\documentclass{article}



\usepackage{graphicx,comment,framed}
\usepackage{roundbox}
\usepackage{fancybox}
\usepackage{tikz}
\usepackage{color}
\usepackage[hidelinks]{hyperref}
\usepackage{framed}
\usepackage{amsthm,amssymb,amsmath}
%\usepackage[colorlinks,linkcolor=blue,citecolor=blue]{hyperref}
\definecolor{shadecolor}{cmyk}{0,0,0,0}
\usepackage{listings}
\usepackage{xepersian}
\usepackage[noend]{algpseudocode}


%--------------------- page settings ----------------------

\settextfont[Scale=1.1]{XB Niloofar}
\setdigitfont[Scale=1.1]{XB Niloofar}
\defpersianfont\sayeh[Scale=1.1]{XB Niloofar}
\addtolength{\textheight}{3.2cm}
\addtolength{\topmargin}{-22mm}
\addtolength{\textwidth}{3cm}
\addtolength{\oddsidemargin}{-1.5cm}


%------------------------ Environments ------------------------------------

\newtheorem{قضیه}{قضیه}
\newtheorem{لم}{لم}
\newtheorem{مشاهده}{مشاهده}
\newtheorem{تعریف}{تعریف}


%-------------------------- Notations ------------------------------------
\renewcommand{\labelitemi}{$\bullet$}
\newcommand{\IR}{\ensuremath{\mathbb{R}}} 
\newcommand{\IZ}{\ensuremath{\mathbb{Z}}} 
\newcommand{\IN}{\ensuremath{\mathbb{N}}} 
\newcommand{\IS}{\ensuremath{\mathbb{S}}} 
\newcommand{\IC}{\ensuremath{\mathbb{C}}} 
\newcommand{\IB}{\ensuremath{\mathbb{B}}} 

\newcommand{\bR}{\mathbb{R}}
\newcommand{\cB}{\mathcal{B}}
\newcommand{\cO}{\mathcal{O}}
\newcommand{\cG}{\mathcal{G}}
\newcommand{\rM}{\mathrm{M}}
\newcommand{\rC}{\mathrm{C}}
\newcommand{\rV}{\mathrm{V}}

\newcommand{\lee}{\leqslant}
\newcommand{\gee}{\geqslant}
\newcommand{\ceil}[1]{{\left\lceil{#1}\right\rceil}}
\newcommand{\floor}[1]{{\left\lfloor{#1}\right\rfloor}}
\newcommand{\prob}[1]{{\mbox{\tt Pr}[#1]}}
\newcommand{\set}[1]{{\{ #1 \}}}
\newcommand{\seq}[1]{{\left< #1 \right>}}
\newcommand{\provided}{\,|\,}
\newcommand{\poly}{\mbox{\rm poly}}
\newcommand{\polylog}{\mbox{\rm \scriptsize polylog}\,}
\newcommand{\comb}[2] {\left(\!\!\begin{array}{c}{#1}\\{#2}\end{array}\!\!\right)}




\newcounter{probcnt}
\newcommand{\مسئله}[1]{\stepcounter{probcnt}{
 	\bf \arabic{probcnt}$\mbox{\bf{.}}$ \ #1}}

\newcommand{\fqed}[1]{\leavevmode\unskip\nobreak\quad\hspace*{\fill}{\ensuremath{#1}}}

\newenvironment{اثبات}
	{\begin{trivlist}\item[\hskip\labelsep{\em اثبات.}]}
	{\fqed{\square}\end{trivlist}}

\newenvironment{حل}
	{\begin{trivlist}\item[\hskip\labelsep{\bf حل.}]}
	{\fqed{\blacktriangleright}\end{trivlist}}

\ifdefined\hidesols
	\newsavebox{\trashcan} % uncomment the following line to hide solutions
	\renewenvironment{حل}{\begin{latin}\begin{lrbox}{\trashcan}}{\end{lrbox}\end{latin}}
\fi


%------------------------- Header -----------------------------

%------------------------- Header -----------------------------

\newcommand{\سربرگ}[3]{
	\parindent=0em
	
	
	\begin{shaded}
		
		\rightline{ 
			\makebox[8em][c]{
				\includegraphics[height=3cm]{aut.png}
		}} \ \
	\\[-3em] 
	\centerline{ مدرس :دکتر امیر مزلقانی } 
	\\[-6.5em]
	\centerline{\large \bf جبرخطی کاربردی }
	\\[0.05em]
	\centerline {\bf نیمسال دوم 98-97}	
	\\[-6.4em]
		\leftline{ 
		\makebox[8em][c]{
			\includegraphics[height=3.2cm]{ceit.jpg}
	}}
	
	
		

		\hrule height .12em
		
		\normalsize
		\vspace{1mm} #1
		\hfill \small  #3
		\vspace{1mm} 
		\hrule height .1em
		
		\vspace{-0.5em} 
		\hfill {\sayeh\large #2} \hfill
	\end{shaded}
	
{\bf توجه !!!}	
\begin{itemize}
	\item 
	سوالات زیر مربوط به فصل اول درس جبر خطی کاربردی با موضوع ((معادلات خطی در جبر خطی )) می باشد که شامل 5 سوال تئوری و 2 سوال عملی است
	\item 
	سوالات را به دقت و مطالعه و به صورت خوانا و مرتب بنویسید
	\item 
	مهلت ارسال این تمرین ساعت 
	\color{red}{23:55 روز شنبه 97/12/18}
	می باشد. 
\end{itemize}
}





\begin{document}
\سربرگ{مجموعه سوالات فصل 4 (فضای برداری)}{}{}

\مسئله{}
در هر مورد مشخص کنید آیا زیر مجموعه ی داده شده یک زیر فضا از فضای برداری مشخص شده می باشد یا خیر.
\begin{enumerate}
	\item 
	تمامی زوج مرتب هایی مثل
	$(x,y)$
	از 
	$\mathbb{R}^2$
	با اعمال جمع برداری و ضرب اسکالر زیر:
	
	$$(x,y)+(x',y')=(x+x',0),c(x,y)=(cx,0)$$
	\item
	$\{(a,b,a+b)|a,b\in \mathbb{R}  \}$
	در فضای برداری 
	$\mathbb{R}^3$.
	\item
	$\{A\in M_n(\mathbb{R})|\det(A)=0\}$
	در
	$M_n(\mathbb{R})$.
	(منظور از 
	$M_n(\mathbb{R})$
	مجموعه تمام ماتریس های 
	$n\times n$
	با درایه هایی از مجموعه اعداد حقیقی است.)
	
	\item 
	$\{p(x)|p(x)=p(-x),p(x)\in\mathbb{P} [x]  \}$
	در فضای برداری 
	$\mathbb{P}[x]$
	(تمامی چند جمله های حداکثر از درجه 
	$n$
	با ضرایب حقیقی را با نماد  
	$\mathbb{P}_n[x]$
	نشان می دهیم و همچنین مجموعه تمام چند جمله ها با ضرایب حقیقی را با  
	$\mathbb{P}[x]$
	نشان می دهیم) 
	\item 
	$\{p(x)|p(x)=ax^2,a\in\mathbb{R}\}$
	در فضای برداری 
	$\mathbb{P}_3[x]$.
	
\end{enumerate}

\مسئله{}
اگر 
$V,W$
فضا ها برداری باشند،
$V\times W$
تمام زوج مرتب هایی به شکل 
$(v,w)$
است که 
$v\in V,w\in W$
و تعریف می کنیم:
$$(v_1,w_1)+(v_2,w_2)=(v_1+v_2,w_1+w_2)$$
و 
$$k(v,w)=(kv,kw),\qquad k\in \mathbb{R}$$
\begin{enumerate}
\item 	
	نشان دهید 
	$V\times W$
	یک فضای برداری است.
	\item  
	نشان دهید اگر بعد 
	$V$
	و
	$W$
	متناهی باشد آنگاه بعد 
	$V\times W$
	نیز متناهی است.
	\item 
	بعد 
	$V\times W$
	را در صورتی که 
	$dimV=m,dimW=n$
	بیابید.
	\item 
	توضیح دهید چرا 
	$\mathbb{R}\times \mathbb{R}^2$
	با 
	$\mathbb{R}^3$
	یکسان است.
	\item 
	پایه ای برای 
	$\mathbb{R}^2\times M_2(\mathbb{R})$
	بیابید.
	\item 
	بعد 
	$M_2(\mathbb{R})\times M_2(\mathbb{R})$
	را بیابید.
\end{enumerate}
\مسئله{}
فرض کنید 
$W_1,W_2$
زیر فضا های فضای برداری 
$V$
باشند، تعریف می کنیم :
$$W_1+W_2=\{w_1+w_2|w_1\in W_1,w_2\in W_2\}$$.

\begin{enumerate}
	\item 
	نشان دهید :
	$$W_1+W_2+\cdots+W_n=span(\bigcup_{i=1}^{n}W_i)$$
	\item 
	نشان دهید 
	$W_1\cap W_2,W_1+W_2$
	زیر فضای 
	$V$
	هستند و همچنین نشان دهید:
	$$W_1\cap W_2 \subseteq W_1\cup W_2\subseteq W_1+W_2$$.
	\item 
	نشان دهید :
	$$dim(W_1+W_2)=dim(W_1)+dim(W_2)-dim(W_1\cap W_2)$$.
	\item 
	نتیجه گیری قسمت 
	$2$
	را با استفاده از دو خط که از مبدا مختصات 
	$xy$
	می گذرند توجیه کنید.
	\item 
	چه زمانی 
	$W_1\cup W_2$
	زیر فضایی از 
	$V$
	است؟
	\item 
	اگر 
	$W_1+W_2$
	کوچکترین زیر فضایی از 
	$V$
	باشد که شامل 
	$W_1\cup W_2$
	است،و اگر 
	$S$
	زیر فضایی از 
	$V$
	باشد که شامل 
	$W_1\cup W_2$
	است آنگاه:
	$W_1+W_2\subseteq S$.
	
\end{enumerate}
\مسئله{}
نشان دهید 
$\{u_1,u_2,\cdots,u_p\}$
در 
$V$
مستقل خطی هستند اگر و فقط اگر بردار های مختصات آن\\

$\{[u_1]_\mathcal{B},[u_2]_\mathcal{B},\cdots,[u_p]_\mathcal{B}\}$
در 
$\mathbb{R}^n$
مستقل خطی باشند.


\مسئله {} فرض کنید 
$S$
یک مجموعه متناهی در 
$V$
باشد به طوری که هر 
$x\in V$
نمایشی یکتا از ترکیب خطی اعضای 
$S$
داشته باشد،نشان دهید 
$S$
پایه 
$V$
است.

\مسئله{}
با در نظر گرفتن 
$T:V\longrightarrow W$
فرض کنید 
$H$
زیر فضایی غیر صفر از 
$V$
باشد،و 
$T(H)$
مجموعه تمام بردار های تصویر شده 
$H$
باشد که زیر فضای 
$W$
است،ثابت کنید
$dim T(H)\leq dim H$.


\مسئله{}ماتریسی را 
(\lr {full rank})
گویند هرگاه رنک آن بزرگترین مقدار ممکن را داشته باشد،ثابت کنید ماتریس 
$m\times n$
که تعداد سطر های آن بیشتر از ستون هایش باشد 
\lr {full rank}
است اگر و فقط اگر ستون های آن مستقل خطی باشند. 

\مسئله {}اگر 
$A$
یک ماتریس 
$m\times n$
با رنک 
$r$
باشد آنگاه می گوییم 
$A=CR$
 یک 
(\lr{rank factorization})
ماتریس 
$A$
است هر گاه
$C$
یک ماتریس 
$m\times r$
با رنک
$r$
و 
$R$
یک ماتریس 
$r\times n$
با رنک 
$r$
باشد.چنین تقسیم بندی همیشه موجود است،با استفاده از این موضوع ثابت کنید برای هر دو ماتریس 
$m\times n$،
$A$
و
$B$،
نشان دهید:
$$rank(A+B)\leq rank A+rank B$$

\مسئله{}(سوال امتیازی)فرض کنید
$V$
 یک فضای برداری متناهی است و 
$V_1,V_2$
زیر فضا هایی از 
$V$
هستند.اگر 


$dim(V_1+V_2)=dim(V_1\cap V_2)+1$
آنگاه 
$V_1+V_2$
یا برابر 
$V_1$
است یا 
$V_2$
و متقابلا 
$V_1\cap V_2 $یا برابر 
$V_1$
است یا 
$V_2$.
هم ارز با جملات قبل اگر 
$V_1,V_2$
زیر مجموعه های هم نباشند،آنگاه:
$$dim(V_1+V_2)\geq dim(V_1\cap V_2)+2$$

\مسئله{}(سوال امتیازی)فرض کنید 
$U$
و
$V$
زیر فضا هایی از 
$R^n$
باشند به طوری که 
$$U=span\{\alpha_1,\alpha_2,\cdots,\alpha_p\}\qquad V=span\{\beta_1,\beta_2,\cdots,\beta_q\}$$
فرض کنید:
$$W=span\{\alpha_i+\beta_j\}\quad i=1,2,\cdots,p \quad j=1,2,\cdots,q$$
اگر 
$dimU=s, dim V=t$
نشان دهید:
$$dimW\leq min\{n,s+t\}$$


\مسئله {} در هر یک از قسمت های زیر ابتدا مختصات بردار داده شده
$(v)$ 
را در هریک از پایه ها بیابید سپس ماتریس انتقال از یک پایه
$(B)$
به پایه
$(C)$
دیگر را محاسبه کنید.
\begin{enumerate}
	\item 
	$$V=\mathbb{P}_3[x]\qquad v=p(x)=8+x+6x^2+9x^3$$
	$$B=\{−2+3x+4x^2-x^3, 3x+5x^2+2x^3, -5x^2-5x^3,
	4 + 4x + 4x^2\}$$
	$$C=\{1 - x^3, 1 + x, x + x^2, x^2 + x^3\}$$
	\item 
	$$V=M_2(\mathbb{R})\qquad v=\begin{bmatrix}
	-3&-2\\
	-1&2
	\end{bmatrix}
	$$
	$$B=\{
	\begin{bmatrix}
	1&0\\
	-1&-2
	\end{bmatrix}
	\begin{bmatrix}
	0&-1\\
	3&0
	\end{bmatrix}
	\begin{bmatrix}
	3&5\\
	0&0
	\end{bmatrix}
	\begin{bmatrix}
	-2&-4\\
	0&0
	\end{bmatrix}
	\}$$
	$$C=\{
	\begin{bmatrix}
	1&1\\
	1&1
	\end{bmatrix}
	\begin{bmatrix}
	1&1\\
	1&0
	\end{bmatrix}
	\begin{bmatrix}
	1&1\\
	0&0
	\end{bmatrix}
	\begin{bmatrix}
	1&0\\
	0&0
	\end{bmatrix}
	\}$$
	\item
	$$V=\mathbb{R}^3\qquad v=(1,7,7)$$
	$$B=\{(-7,4,4),(4,2,-1),(-7,5,0)\}$$
	$$C={(1,1,0),(0,1,1),(3,-1,-1)}$$
\end{enumerate}

\end{document}


